\chapter{Computationele chemie: Hartree-Fock en DFT}\label{ch:b3:computational-chemistry}

Het heliumhydride-ion, \ce{HeH+}, geldt als het eerste molecuul dat ontstond
toen het jonge heelal afkoelde: een heliumatoom dat een proton vasthoudt.
Het werd in 1925 gemaakt in een gasontladingsbuis in het laboratorium, maar
er werd decennialang tevergeefs naar gezocht in de ruimte, tot zijn
rotatielijn in 2019 werd waargenomen in een
planetaire nevel. Zijn bindingslengte, zijn vibratiefrequenties en de
golflengten van zijn lijnen waren toen al lang bekend --- uit
berekeningen. Een computerprogramma dat de \omterm{thm:b3:quantum-model-systems:schrodinger}{schrödingervergelijking} van de
elektronen bij benadering oplost, kan de vorm, de energie
en het spectrum voorspellen van een molecuul dat niemand ooit in een fles
heeft gestopt. Dit hoofdstuk legt uit hoe zulke programma's werken, van de
scheiding van kernen en elektronen tot de hartree-focktheorie en de
\omterm{def:b3:computational-chemistry:dft}{dichtheidsfunctionaaltheorie}, en waarvoor hun getallen wel en niet te
vertrouwen zijn.

\begin{recall}
\Cref{ch:b3:quantum-model-systems} gaf de \omterm{def:b3:quantum-model-systems:schrodinger}{hamiltonoperator}, de
\omterm{thm:b3:quantum-model-systems:schrodinger}{schrödingervergelijking} en het waterstofatoom; \cref{ch:b3:many-electron-atoms} de
\omterm{def:b3:many-electron-atoms:slater}{slaterdeterminant} en de \omterm{def:b3:many-electron-atoms:exchange}{uitwisselingsintegraal}. Het deel van jaar~2
bouwde moleculaire orbitalen als lineaire combinaties van atoomorbitalen
(LCAO), met overlap-, coulomb- en resonantie-integralen en een seculiere
determinant, en loste de hückelmethode op; het gebruikte ook
foto-elektronspectra om orbitaalenergieën zichtbaar te maken.
\end{recall}

\begin{omfigure}
\includegraphics[width=0.6\linewidth]{images/book4/ai/b3-03-cluster.jpg}

{\small Een rekencluster: rijen servers waarop de meeste
kwantumchemische berekeningen tegenwoordig draaien, vaak uren of dagen
per molecuul.}
\end{omfigure}

\section{Born-Oppenheimer en het potentiaalenergieoppervlak}

De \omterm{def:b3:quantum-model-systems:schrodinger}{hamiltonoperator} van een molecuul bevat de kinetische energieën van
zijn kernen en elektronen en al hun coulombwisselwerkingen. Kernen zijn
minstens 1836 keer zwaarder dan elektronen: bij dezelfde krachten
bewegen ze veel trager, en de elektronen passen zich vrijwel onmiddellijk
aan elke positie van de kernen aan.

\begin{definition}[Born-oppenheimerbenadering]\label{def:b3:computational-chemistry:born-oppenheimer}
In de \emph{born-oppenheimerbenadering}\index{born-oppenheimerbenadering}
wordt de elektronische \omterm{thm:b3:quantum-model-systems:schrodinger}{schrödingervergelijking} opgelost voor kernen die
vastgehouden worden op de posities $\mathbf R$,
\[
\hat H_{\mathrm{el}}\,\psi_{\mathrm{el}}(\mathbf r;\mathbf R) =
E_{\mathrm{el}}(\mathbf R)\,\psi_{\mathrm{el}}(\mathbf r;\mathbf R),
\]
en bewegen de kernen vervolgens in de potentiaal $E_{\mathrm{el}}(\mathbf R) +
V_{\mathrm{nn}}(\mathbf R)$, de elektronische energie plus de afstoting
van de kernen.
\end{definition}

\begin{definition}[Potentiaalenergieoppervlak, stationair punt, geometrieoptimalisatie]\label{def:b3:computational-chemistry:pes}
Het \emph{potentiaalenergieoppervlak}\index{potentiaalenergieoppervlak} (PES) van
een molecuul is de functie $U(\mathbf R) = E_{\mathrm{el}}(\mathbf R) +
V_{\mathrm{nn}}(\mathbf R)$ van zijn kerncoördinaten. Een \emph{stationair
punt}\index{stationair punt} is een punt waar alle eerste afgeleiden van $U$
nul zijn; een minimum is een
\emph{evenwichtsgeometrie}\index{evenwichtsgeometrie}. Het zoeken van zo'n
minimum vanuit een beginstructuur, door de
krachten $-\nabla U$ te volgen, is een
\emph{geometrieoptimalisatie}\index{geometrieoptimalisatie}.
\end{definition}

Voor een diatomisch molecuul is het PES een kromme $U(R)$; voor een
triatomisch molecuul een functie van drie coördinaten; voor $N$ atomen van
$3N - 6$. De massa van de
kernen komt niet voor in $\hat H_{\mathrm{el}}$: \ce{H2}, \ce{HD} en \ce{D2} delen
hetzelfde PES, dezelfde bindingslengte en \omterm{def:b3:quantum-model-systems:oscillator}{krachtconstante}, en verschillen
alleen in de manier waarop de kernen erop bewegen (\omterm{def:b3:quantum-model-systems:oscillator}{nulpuntsenergie},
\cref{ch:b3:quantum-model-systems}).

\begin{proposition}[Frequenties uit de hessematrix]\label{prop:b3:computational-chemistry:hessian}
Nabij een minimum is $U \approx U_0 + \frac12\sum_{ij}H_{ij}\,\delta q_i\,\delta q_j$, met
de hessematrix $H_{ij} = \partial^2U/\partial q_i\partial q_j$. In massagewogen
coördinaten $\delta q_i\sqrt{m_i}$ geven de \omterm{def:b3:quantum-model-systems:operator}{eigenwaarden} $\lambda_k$ van de
massagewogen hessematrix de harmonische hoekfrequenties $\omega_k =
\sqrt{\lambda_k}$ van de normaaltrillingen. In een minimum zijn alle $\lambda_k$
positief (zes ervan, of vijf voor een lineair molecuul, zijn nul:
translaties en rotaties); in een zadelpunt van de eerste orde is er precies
één negatief.
\end{proposition}

\begin{proof}[Gedeeltelijk bewijs]
Voor een diatomisch molecuul is $U \approx U_0 + \frac12k(R - R_e)^2$ met $k = U''(R_e)$, en
de oscillator van \cref{ch:b3:quantum-model-systems} geeft $\omega =
\sqrt{k/\mu}$. In het algemeen worden de klassieke vergelijkingen $m_i\ddot q_i =
-\sum_jH_{ij}q_j$, in massagewogen coördinaten $x_i = \sqrt{m_i}q_i$,
$\ddot{\mathbf x} = -\tilde H\mathbf x$ met de symmetrische matrix $\tilde H_{ij} =
H_{ij}/\sqrt{m_im_j}$; haar eigenvectoren oscilleren onafhankelijk met
$\sqrt{\lambda_k}$. Een negatieve $\lambda_k$ geeft een imaginaire frequentie:
in die richting daalt de energie, een maximum. De normaaltrillingen van
\cref{ch:b3:group-theory-applied} zijn deze eigenvectoren.
\end{proof}

\begin{omfigure}
\begin{tikzpicture}[font=\small, scale=0.95]
  \fill[omDef!8] (0,0) ellipse (1.9 and 1.15);
  \foreach \s in {1.0,0.75,0.5,0.25}{\draw[omDef!70] (0,0) ellipse ({1.9*\s} and {1.15*\s});}
  \fill[omThm!8] (6,0.4) ellipse (1.6 and 1.0);
  \foreach \s in {1.0,0.75,0.5,0.25}{\draw[omThm!70] (6,0.4) ellipse ({1.6*\s} and {1.0*\s});}
  \draw[gray] (2.1,-0.15) .. controls (2.9,0.55) and (3.4,0.55) .. (4.3,0.35);
  \draw[gray] (2.2,0.9) .. controls (3.0,0.35) and (3.5,0.35) .. (4.1,1.2);
  \node[omsaddle] at (3.15,0.42) {};
  \draw[ompath] (0,0) .. controls (1.8,0.6) and (2.4,0.5) .. (3.15,0.42) .. controls (4.0,0.35) and (4.8,0.4) .. (6,0.4);
  \node[anchor=north] at (0,-1.25) {minimum A (beginstoffen)};
  \node[anchor=north] at (6,-0.7) {minimum B (producten)};
  \node[anchor=south] at (3.15,1.25) {zadelpunt};
  \draw[->, gray] (3.15,1.25) -- (3.15,0.6);
  \fill (0,0) circle (1.2pt) (6,0.4) circle (1.2pt);
  \draw[->] (-2.3,-2.2) -- (8.2,-2.2) node[right] {$q_1$};
  \draw[->] (-2.3,-2.2) -- (-2.3,1.8) node[above] {$q_2$};
\end{tikzpicture}

{\small Een \omterm{def:b3:computational-chemistry:pes}{potentiaalenergieoppervlak} van bovenaf gezien, als
hoogtelijnenkaart boven twee kerncoördinaten (schematisch). Twee minima
zijn via een zadelpunt (kruisje) verbonden door het pad van laagste
energie (gestreept): het pad van minimale energie dat in
\cref{ch:b3:rate-theories} bestudeerd wordt.}
\end{omfigure}

\section{Het variatieprincipe}

De elektronische \omterm{thm:b3:quantum-model-systems:schrodinger}{schrödingervergelijking} kan voor geen enkel molecuul met
meer dan één elektron exact opgelost worden. Benaderingen worden
beoordeeld aan de hand van één stelling.

\begin{theorem}[Het variatieprincipe]\label{thm:b3:computational-chemistry:variational}
Voor elke normeerbare proeffunctie $\phi$ die aan de randvoorwaarden van
het probleem voldoet, geldt
\[
E[\phi] = \frac{\langle\phi|\hat H|\phi\rangle}{\langle\phi|\phi\rangle} \ge E_0,
\]
waarin $E_0$ de laagste \omterm{def:b3:quantum-model-systems:operator}{eigenwaarde} van $\hat H$ is; gelijkheid geldt alleen
als $\phi$ een \omterm{def:b3:quantum-model-systems:operator}{eigenfunctie} van de grondtoestand is. Dit is het
\emph{variatieprincipe}\index{variatieprincipe}.
\end{theorem}

\begin{proof}
Ontwikkel $\phi = \sum_nc_n\psi_n$ naar de orthonormale \omterm{def:b3:quantum-model-systems:operator}{eigenfuncties} van $\hat H$
(\cref{thm:b3:quantum-model-systems:hermitian-real}). Dan is $\langle\phi|\hat
H|\phi\rangle = \sum_n|c_n|^2E_n \ge E_0\sum_n|c_n|^2 = E_0\langle\phi|\phi\rangle$,
omdat elke $E_n \ge E_0$; gelijkheid vereist $c_n = 0$ telkens wanneer $E_n > E_0$.
\end{proof}

Hoe lager de energie van een proeffunctie, hoe beter ze is; elke
aanpasbare parameter wordt vastgelegd door $E$ te minimaliseren.

\begin{proposition}[Helium met een afgeschermde lading]\label{prop:b3:computational-chemistry:helium-zeta}
Voor helium geeft de proeffunctie $\phi = \eu^{-\zeta(r_1 + r_2)}$ (in atomaire
eenheden: lengten in $a_0$, energieën in $E_h$) $E(\zeta) = \zeta^2 - \frac{27}{8}\zeta$,
minimaal voor $\zeta = \frac{27}{16}$, met $E = -(\frac{27}{16})^2 E_h =
\qty{-2.8477}{\hartree}$.
\end{proposition}

\begin{proof}
Voor één elektron in $\eu^{-\zeta r}$ is $\langle T\rangle = \zeta^2/2$ en
$\langle 1/r\rangle = \zeta$; de afstoting van twee elektronen in dezelfde
$1s$-functie is $\frac58\zeta$ (een standaardintegraal, hier aangenomen).
Dus $E = 2\cdot\frac12
\zeta^2 - 2\cdot2\zeta + \frac58\zeta = \zeta^2 - \frac{27}{8}\zeta$, en $\dd E/\dd
\zeta = 0$ geeft $\zeta = 27/16 = 1.6875$.
\end{proof}

\begin{example}[Hoe goed is dat?]\label{ex:b3:computational-chemistry:helium}
% ledger: ie:He, ie4:He2, const:Eh, const:eV
De totale energie van helium is min de som van zijn twee
ionisatie-energieën,
$-(24.5874 + 54.4178)\,\unit{eV} = \qty{-2.9034}{\hartree}$. De functie met
één parameter ligt \qty{0.056}{\hartree} (\qty{1.5}{eV}) te hoog: elk
elektron ziet een kern die door het andere afgeschermd is tot
$Z_{\mathrm{eff}} = 1.69$, zoals de regels van Slater
suggereren, maar de elektronen ontwijken elkaar ook op een manier die geen
product van twee functies kan beschrijven.
\end{example}

\begin{proposition}[Lineaire variatie]\label{prop:b3:computational-chemistry:linear-variation}
Voor een proeffunctie $\phi = \sum_{i=1}^n c_i\chi_i$ opgebouwd uit $n$ vaste functies,
met $H_{ij} = \langle\chi_i|\hat H|\chi_j\rangle$ en $S_{ij} = \langle\chi_i|\chi_j
\rangle$, zijn de stationaire waarden van $E$ de wortels van $\det(H - ES) = 0$, en
de laagste wortel is een bovengrens voor $E_0$.
\end{proposition}

\begin{proof}
$E\sum_{ij}c_ic_jS_{ij} = \sum_{ij}c_ic_jH_{ij}$ (reële coëfficiënten).
Differentieer naar $c_k$ en stel $\partial E/\partial c_k = 0$: $\sum_j(H_{kj} -
ES_{kj})c_j = 0$ voor elke $k$. Een oplossing die niet nul is, vereist een
determinant die nul is. De laagste wortel is $E[\phi]$ voor zijn
eigenvector, dus $\ge E_0$.
\end{proof}

De hückelmethode uit het deel van jaar~2 is deze propositie met
$p$-orbitalen en empirische $H_{ij}$. Elke methode hieronder is hetzelfde
idee met betere functies en betere matrixelementen.

\section{Basissets}

\begin{definition}[Basisset, minimale basisset, slater- en gaussorbitalen]\label{def:b3:computational-chemistry:basis}
Een \emph{basisset}\index{basisset} is de verzameling vaste functies $\chi_i$,
gecentreerd op de atomen, waaruit moleculaire orbitalen opgebouwd worden.
Een \emph{minimale basisset}\index{minimale basisset} heeft één functie per
bezette atoomorbitaal (één $1s$ voor H, vijf voor C). Een
\emph{slaterorbitaal}\index{slaterorbitaal} (STO) heeft de radiale vorm
$r^{n-1}\eu^{-\zeta r}$; een \emph{gaussorbitaal}\index{gaussorbitaal}
(GTO) heeft de vorm $\eu^{-\alpha r^2}$ maal een polynoom in $x$, $y$,
$z$. Een \emph{gecontraheerde gaussfunctie}\index{gecontraheerde gaussfunctie}
is een vaste combinatie $\sum_kd_k\,\eu^{-\alpha_kr^2}$ van primitieve gaussfuncties.
\end{definition}

Slaterfuncties hebben de juiste vorm (een knik in de kern, een
exponentiële staart), maar hun tweeëlektronintegralen over meerdere
centra zijn zeer duur; gaussfuncties hebben de verkeerde vorm, maar het
product van twee gaussfuncties op verschillende centra is een gaussfunctie
op een punt daartussen, zodat elke integraal herleid wordt tot een
gesloten formule. Contracties combineren het beste van beide: in de basis
STO-3G wordt elke slaterfunctie vervangen door een vaste som van drie
gaussfuncties die erop zijn aangepast.

\begin{example}[De STO-3G-functie van waterstof]\label{ex:b3:computational-chemistry:sto3g}
% ledger: bse:H.sto3g.a1, bse:H.sto3g.a2, bse:H.sto3g.a3, bse:sto3g.c1, bse:sto3g.c2, bse:sto3g.c3
De $1s$-functie van waterstof in STO-3G is $0.15433\,g(3.42525) + 0.53533\,
g(0.62391) + 0.44463\,g(0.16886)$, waarin $g(\alpha)$ de genormeerde gaussfunctie
$(2\alpha/\pi)^{3/4}\eu^{-\alpha r^2}$ is (exponenten in $a_0^{-2}$). Ze bootst een
slaterfunctie met exponent $\zeta = 1.24$ na --- een waterstofatoom dat een
beetje samengedrukt is, zoals in moleculen. Haar overlap met de exacte
$1s$-functie met dezelfde exponent is $0.9998$; alleen de knik in de kern
en de verre staart worden gemist (figuur hieronder).
\end{example}

\begin{definition}[Gesplitste-valentiebasisset, polarisatie- en diffuse functies]\label{def:b3:computational-chemistry:split-valence}
Een \emph{gesplitste-valentiebasisset}\index{gesplitste-valentiebasisset}
beschrijft elke valentieorbitaal met twee (of meer) functies van
verschillende grootte, zodat de orbitaal in een molecuul kan groeien of
krimpen. Een \emph{polarisatiefunctie}\index{polarisatiefunctie} heeft een
hogere $l$ dan de bezette
orbitalen ($p$ op H, $d$ op C) en laat de dichtheid van het atoom
wegschuiven; een \emph{diffuse functie}\index{diffuse functie} heeft een
kleine exponent en beschrijft anionen en zwakke wisselwerkingen ver van de
kernen.
\end{definition}

Een naam als 6-31G(d) zegt: zes primitieven voor elke kernfunctie, de
valentie gesplitst in een contractie van drie en één enkele primitieve, en
$d$-polarisatiefuncties op de zware atomen. De resultaten convergeren naar
de \emph{basissetlimiet} naarmate de basis groeit; een berekening is maar
zo goed als haar basis.

\begin{omfigure}
\begin{tikzpicture}
\begin{axis}[width=0.46\linewidth, height=5.4cm, xmin=0, xmax=5, ymin=0,
  ymax=0.6, axis lines=left, xlabel={$r$ / $a_0$}, ylabel={$\phi(r)$},
  legend cell align=left, legend style={font=\small, draw=none},
  tick label style={font=\small}, label style={font=\small},
  title={\small STO-3G tegenover $\eu^{-r}$ ($\zeta = 1$)}]
\addplot[black, thick] table[x=r, y=slater] {figdata/out/bachelor-3-minimal-basis-hf-sto.dat};
\addlegendentry{Slater}
\addplot[omDef, thick, dashed] table[x=r, y=g3] {figdata/out/bachelor-3-minimal-basis-hf-sto.dat};
\addlegendentry{STO-3G}
\addplot[omThm, thick, dotted] table[x=r, y=g1] {figdata/out/bachelor-3-minimal-basis-hf-sto.dat};
\addlegendentry{één gaussfunctie}
\end{axis}
\end{tikzpicture}\hfill
\begin{tikzpicture}
\begin{axis}[width=0.46\linewidth, height=5.4cm, xmin=0.5, xmax=8.5,
  ymin=-2.86, ymax=-2.78, axis lines=left, xlabel={SCF-cyclus},
  ylabel={$E$ / $E_h$}, ytick={-2.85,-2.83,-2.81,-2.79},
  yticklabel style={/pgf/number format/fixed, /pgf/number format/precision=2},
  tick label style={font=\small}, label style={font=\small},
  title={\small SCF van \ce{HeH+}}]
\addplot[omDef, thick, mark=*] table[x=it, y=E] {figdata/out/bachelor-3-minimal-basis-hf-heh-scf.dat};
\addplot[gray, dashed, domain=0.5:8.5] {-2.84184};
\end{axis}
\end{tikzpicture}

{\small Links: de STO-3G-contractie (gestreept) geeft de $1s$-slaterfunctie
nauwkeurig weer, behalve in de kern; één enkele gaussfunctie (gestippeld)
doet dat niet.
Rechts: de energie van \ce{HeH+} bij \qty{1.4632}{\bohr} tijdens de
iteraties van het zelfconsistente veld (STO-3G); ze convergeert van bovenaf
naar \qty{-2.8418}{\hartree} in enkele cycli, en elke iteratie gehoorzaamt
aan het \omterm{thm:b3:computational-chemistry:variational}{variatieprincipe}.}
\end{omfigure}

\section{Hartree-focktheorie}

De eenvoudigste \omterm{def:b3:many-electron-atoms:indistinguishable}{antisymmetrische golffunctie} van $N$ elektronen is één
\omterm{def:b3:many-electron-atoms:slater}{slaterdeterminant}. De hartree-focktheorie zoekt de beste.

\begin{definition}[Hartree-fockmethode, fockoperator, zelfconsistent veld]\label{def:b3:computational-chemistry:hartree-fock}
De \emph{hartree-fockmethode}\index{hartree-fockmethode} benadert de
grondtoestand door de \omterm{def:b3:many-electron-atoms:slater}{slaterdeterminant} met de laagste energie. Haar
orbitalen zijn \omterm{def:b3:quantum-model-systems:operator}{eigenfuncties} van de \emph{fockoperator}\index{fockoperator}
\[
\hat f(1) = \hat h(1) + \sum_b\big[2\hat J_b(1) - \hat K_b(1)\big],
\]
in het geval van gesloten schillen: $\hat h$ is de kinetische energie en de
aantrekking door de kernen van één elektron, $\hat J_b$ de afstoting door
de ladingswolk van orbitaal $b$, en $\hat K_b$ de uitwisselingsoperator,
waarvan de \omterm{def:b3:quantum-model-systems:hermitian}{verwachtingswaarden} de uitwisselingsintegralen van
\cref{ch:b3:many-electron-atoms} zijn. Omdat $\hat f$ afhangt van de
orbitalen die hij bepaalt, worden de vergelijkingen iteratief opgelost tot
de orbitalen niet meer veranderen: een \emph{zelfconsistent
veld}\index{zelfconsistent veld} (SCF).
\end{definition}

\begin{theorem}[Vergelijkingen van Roothaan-Hall]\label{thm:b3:computational-chemistry:roothaan}
Als elke orbitaal in een basis geschreven wordt, $\phi_a = \sum_\mu C_{\mu a}\chi_\mu$,
worden de hartree-fockvergelijkingen voor gesloten schillen het
matrixeigenwaardeprobleem
\[
\mathbf F\mathbf C = \mathbf S\mathbf C\boldsymbol\varepsilon, \qquad F_{\mu\nu} =
h_{\mu\nu} + \sum_{\lambda\sigma}P_{\lambda\sigma}\big[(\mu\nu|\sigma\lambda) -
\tfrac12(\mu\lambda|\sigma\nu)\big],
\]
met de dichtheidsmatrix $P_{\lambda\sigma} = 2\sum_a^{\mathrm{occ}}C_{\lambda a}C_{\sigma
a}$ en de tweeëlektronintegralen $(\mu\nu|\lambda\sigma) = \iint\chi_\mu(1)\chi_\nu(1)
r_{12}^{-1}\chi_\lambda(2)\chi_\sigma(2)$ (atomaire eenheden).
\end{theorem}

\begin{proof}[Gedeeltelijk bewijs]
De energie van de determinant is $E = \sum_{\mu\nu}P_{\mu\nu}h_{\mu\nu} + \frac12
\sum P_{\mu\nu}P_{\lambda\sigma}[(\mu\nu|\sigma\lambda) - \frac12(\mu\lambda|\sigma\nu)]$, een
kwadratische functie van de coëfficiënten. Minimaliseer haar onder de
voorwaarden dat de orbitalen orthonormaal blijven, $\mathbf C^{\mathsf T}\mathbf S\mathbf C = \mathbf 1$,
met een lagrangemultiplicator voor elke voorwaarde: de
stationariteitsvoorwaarden zijn $\mathbf F\mathbf C = \mathbf S\mathbf C\boldsymbol\varepsilon$, waarbij de
multiplicatoren een matrix vormen die diagonaal gemaakt kan worden door
een rotatie van de bezette orbitalen onder elkaar. De algebra wordt hier
aangenomen.
\end{proof}

\begin{method}[De SCF-procedure]\label{met:b3:computational-chemistry:scf}
\begin{enumerate}
  \item Bereken eenmalig de integralen $S_{\mu\nu}$, $h_{\mu\nu}$ en $(\mu\nu|\lambda\sigma)$.
  \item Schat de dichtheidsmatrix (bijvoorbeeld uit de eigenvectoren van
        $\mathbf h$ alleen, de ``core''-schatting).
  \item Stel $\mathbf F$ op uit $\mathbf P$; los $\mathbf F\mathbf C = \mathbf S\mathbf C
        \boldsymbol\varepsilon$ op (orthogonaliseer met $\mathbf S^{-1/2}$ en
        diagonaliseer daarna).
  \item Vul de laagste orbitalen, met twee elektronen elk; vorm de nieuwe
        $\mathbf P$ en de energie.
  \item Herhaal vanaf stap 3 tot de energie en $\mathbf P$ minder dan een
        drempelwaarde veranderen.
\end{enumerate}
\end{method}

\begin{omfigure}
\begin{tikzpicture}[font=\small, node distance=4mm,
  box/.style={draw=black!60, rounded corners=2pt, fill=omDef!6, align=center,
  text width=4.4cm, inner sep=4pt}]
  \node[box] (a) {integralen $S$, $h$, $(\mu\nu|\lambda\sigma)$};
  \node[box, below=of a] (b) {schat $\mathbf P$};
  \node[box, below=of b] (c) {stel $\mathbf F(\mathbf P)$ op};
  \node[box, below=of c] (d) {los $\mathbf F\mathbf C = \mathbf S\mathbf C\boldsymbol\varepsilon$ op};
  \node[box, below=of d] (e) {nieuwe $\mathbf P$, energie $E$};
  \node[diamond, aspect=2.4, draw=black!60, fill=omThm!6, below=of e, inner sep=1pt] (f) {geconvergeerd?};
  \node[box, below=of f, fill=omProp!8] (g) {orbitalen, $\varepsilon$, $E$};
  \draw[->] (a) -- (b); \draw[->] (b) -- (c); \draw[->] (c) -- (d);
  \draw[->] (d) -- (e); \draw[->] (e) -- (f);
  \draw[->] (f) -- node[left] {ja} (g);
  \draw[->] (f.east) -- ++(1.6,0) |- node[pos=0.25, right] {nee} (c.east);
\end{tikzpicture}

{\small De lus van het zelfconsistente veld in een hartree-fockberekening.}
\end{omfigure}

\begin{example}[\ce{H2} in de minimale basis]\label{ex:b3:computational-chemistry:h2}
% ledger: re:H2, janaf:H.dfH0, diat:H2.we, diat:H2.wexe
Met één STO-3G-functie per atoom ligt de bindende orbitaal vast door de
symmetrie,
$\sigma_g \propto \chi_A + \chi_B$, en heeft de SCF maar één stap nodig. Bij $R =
\qty{1.4}{\bohr}$ is de overlap $S_{AB} = 0.6593$ en de energie
\qty{-1.1167}{\hartree}; het minimum ligt bij \qty{1.346}{\bohr} (\qty{71.2}{pm},
tegenover de gemeten \qty{74.1}{pm}) en \qty{-1.1175}{\hartree}. De figuur
hieronder vergelijkt de kromme met de echte, opgebouwd uit gemeten
spectroscopische gegevens en de dissociatie-energie (een morsekromme,
\cref{ch:b3:rovibrational-spectroscopy}).
\end{example}

\begin{omfigure}
\begin{tikzpicture}
\begin{axis}[width=0.7\linewidth, height=6cm, xmin=0.8, xmax=10, ymin=-1.2,
  ymax=-0.45, axis lines=left, xlabel={$R$ / $a_0$}, ylabel={$E$ / $E_h$},
  legend cell align=left, legend style={font=\small, at={(0.04,0.98)}, anchor=north west, draw=none},
  tick label style={font=\small}, label style={font=\small}]
\addplot[omDef, thick] table[x=R, y=rhf] {figdata/out/bachelor-3-minimal-basis-hf-h2.dat};
\addlegendentry{RHF/STO-3G}
\addplot[black, thick, dashed] table[x=R, y=morse] {figdata/out/bachelor-3-minimal-basis-hf-h2.dat};
\addlegendentry{gemeten (morsevorm)}
\addplot[gray, dotted, domain=0.8:10] {-1.0};
\addplot[omThm, dotted, domain=0.8:10] {-0.93316};
\node[font=\scriptsize, anchor=north] at (axis cs:8,-1.005) {twee exacte H-atomen};
\node[font=\scriptsize, anchor=south, omThm] at (axis cs:8,-0.930) {twee STO-3G-H-atomen};
\end{axis}
\end{tikzpicture}

{\small De energie van \ce{H2} als functie van de kernafstand: beperkte
Hartree-Fock in de basis STO-3G (doorgetrokken) en de echte kromme
(gestreept, morsevorm opgebouwd uit de gemeten $D_0$, $\omega_e$, $\omega_ex_e$, $r_e$). Nabij het
minimum is RHF goed; bij grote $R$ stijgt de energie ver boven die van
twee waterstofatomen.}
\end{omfigure}

\begin{proposition}[Waarom beperkte Hartree-Fock faalt bij dissociatie]\label{prop:b3:computational-chemistry:rhf-dissociation}
In de minimale basis is de RHF-golffunctie van \ce{H2}, op de normering na,
$\sigma_g(1)\sigma_g(2) \propto \chi_A(1)\chi_B(2) + \chi_B(1)\chi_A(2)
+ \chi_A(1)\chi_A(2) + \chi_B(1)\chi_B(2)$: op elke afstand half covalent en half
ionair (\ce{H+ H-}). Bij grote $R$ nadert haar energie het gemiddelde van
twee neutrale atomen en een ionenpaar, niet de energie van twee atomen.
\end{proposition}

\begin{proof}
Werk $(\chi_A + \chi_B)(1)(\chi_A + \chi_B)(2)$ uit: de twee kruistermen
plaatsen één elektron op elk atoom, de twee kwadratische termen beide op
hetzelfde atoom, met gelijke gewichten, wat $R$ ook is. De determinant kan
deze gewichten niet veranderen.
\end{proof}

\begin{definition}[Elektronencorrelatie, correlatie-energie]\label{def:b3:computational-chemistry:correlation}
\emph{Elektronencorrelatie}\index{elektronencorrelatie} is het deel van het
wederzijds ontwijken van de elektronen dat niet door één enkele
determinant beschreven wordt (waarin elektronen met tegengestelde spin
onafhankelijk van elkaar bewegen in elkaars gemiddelde veld). De
\emph{correlatie-energie}\index{correlatie-energie} is de exacte
energie min de hartree-fockenergie in een volledige basis; ze is altijd
negatief.
\end{definition}

De \omterm{def:b3:computational-chemistry:correlation}{correlatie-energie} bedraagt ongeveer \qty{-0.04}{\hartree} per
elektronenpaar, één procent van een totale energie, maar vergelijkbaar met
de energie van een reactie: methoden die haar toevoegen
(configuratie-interactie, storingstheorie, coupled cluster) kosten veel
meer dan Hartree-Fock, waarvan de rekenkost al ruwweg met de vierde macht
van het aantal basisfuncties groeit.

\begin{theorem}[Stelling van Koopmans]\label{thm:b3:computational-chemistry:koopmans}
In de hartree-focktheorie is de energie die nodig is om een elektron uit
de bezette orbitaal $a$ te verwijderen, terwijl alle andere orbitalen
bevroren blijven, gelijk aan $-\varepsilon_a$. De
gemeten ionisatie-energieën zijn dus bij benadering de tegengestelden van
de orbitaalenergieën: dit is de
\emph{stelling van Koopmans}\index{stelling van Koopmans}.
\end{theorem}

\begin{proof}
De energie van een determinant met gesloten schillen is $\sum_a2h_{aa} + \sum_{ab}(2J_{ab} -
K_{ab})$, en $\varepsilon_a = h_{aa} + \sum_b(2J_{ab} - K_{ab})$. Verwijder één
elektron uit $a$ zonder de orbitalen te veranderen: de verloren termen zijn
$h_{aa}$, de wisselwerkingen van dat elektron met alle andere,
$\sum_b(2J_{ab} - K_{ab}) -
J_{aa}$, plus $J_{aa}$ met zijn vroegere partner --- in totaal precies
$\varepsilon_a$. Dus $E^+ - E = -\varepsilon_a$.
\end{proof}

\begin{example}[\ce{H2} ioniseren]\label{ex:b3:computational-chemistry:koopmans}
% ledger: ie4:H2
Bij \qty{1.4}{\bohr} is de STO-3G-orbitaalenergie van $\sigma_g$ gelijk aan
\qty{-0.5782}{\hartree}: Koopmans voorspelt \qty{15.73}{eV}, tegenover de
gemeten \qty{15.43}{eV}. De bevroren orbitalen overschatten de energie van
het ion (dat zou relaxeren), de ontbrekende correlatie werkt in de andere
richting; bij valentie-ionisaties heffen de twee fouten elkaar grotendeels
op, en daarom kunnen de foto-elektronspectra uit het deel van jaar~2 met
orbitaaldiagrammen gelezen worden.
\end{example}

\section{Dichtheidsfunctionaaltheorie en wat een berekening vertelt}

De golffunctie van $N$ elektronen hangt af van $3N$ coördinaten; de
\omterm{def:b3:computational-chemistry:dft}{elektronendichtheid} van drie.

\begin{definition}[Elektronendichtheid, dichtheidsfunctionaaltheorie]\label{def:b3:computational-chemistry:dft}
De \emph{elektronendichtheid}\index{elektronendichtheid} $\rho(\mathbf r)$ is het
aantal elektronen per volume-eenheid in $\mathbf r$, gesommeerd over alle
elektronen:
$\int\rho\,\dd\tau = N$. De
\emph{dichtheidsfunctionaaltheorie}\index{dichtheidsfunctionaaltheorie}
(DFT) berekent de energie van de grondtoestand als een functionaal van
$\rho$. In de praktijk wordt $\rho = \sum_a|\phi_a|^2$ opgebouwd uit
\emph{kohn-shamorbitalen}\index{kohn-shamorbitalen}, de orbitalen van
fictieve onafhankelijke elektronen met dezelfde dichtheid als de echte, en
worden alle effecten van uitwisseling en correlatie samengebracht in een
\emph{uitwisselings-correlatiefunctionaal}\index{uitwisselings-correlatiefunctionaal}
$E_{\mathrm{xc}}[\rho]$.
\end{definition}

\begin{theorem}[Hohenberg-Kohn]\label{thm:b3:computational-chemistry:hohenberg-kohn}
De dichtheid van de grondtoestand van een systeem van elektronen bepaalt
de externe potentiaal (de kernen) op een constante na, en dus de
\omterm{def:b3:quantum-model-systems:schrodinger}{hamiltonoperator} en elke eigenschap van de grondtoestand; en de
energiefunctionaal van de dichtheid is minimaal bij de echte dichtheid.
\end{theorem}
\admitted

Het bewijs hoort thuis in meer gevorderde cursussen. Het garandeert dat er
een exacte functionaal bestaat, niet hoe die eruitziet: elke praktische
functionaal is een benadering, gekozen uit een hiërarchie --- lokale
dichtheid, gradiëntcorrecties, hybriden die er wat hartree-fockuitwisseling
bij mengen --- en geijkt op referentiegegevens. De kohn-shamvergelijkingen
worden opgelost met dezelfde SCF-lus als Hartree-Fock, tegen een
vergelijkbare kost, maar bevatten correlatie; daarom werd DFT het
werkpaard van de chemie.

\begin{method}[Een berekening lezen]\label{met:b3:computational-chemistry:read-calculation}
\begin{enumerate}
  \item Controleer dat de SCF geconvergeerd is en dat de
        \omterm{def:b3:computational-chemistry:pes}{geometrieoptimalisatie} eindigde met krachten onder de drempelwaarde.
  \item Voer een frequentieberekening uit op hetzelfde niveau: alleen reële
        frequenties betekent een minimum; één imaginaire frequentie, een
        overgangsstructuur, waarvan de normaaltrilling de beweging over de
        barrière is.
  \item Vergelijk gelijksoortige grootheden: energieverschillen tussen
        structuren die op één niveau berekend zijn, nooit absolute
        energieën van verschillende niveaus.
  \item Vergelijk met het experiment en ken daarbij de typische fouten van
        het niveau: harmonische frequenties worden meestal geschaald met
        een factor 0.9 tot 0.97; bindingslengten kloppen tot op ongeveer
        1 pm met een hybride functionaal en een gepolariseerde basis;
        reactie-energieën in het beste geval tot op enkele \unit{kJ/mol}.
\end{enumerate}
\end{method}

\begin{method}[Een theorieniveau kiezen]\label{met:b3:computational-chemistry:choose-level}
Geometrieën en frequenties van organische moleculen: een hybride
functionaal met een gepolariseerde gesplitste-valentiebasis. Zwakke
wisselwerkingen: voeg een dispersiecorrectie en \omterm{def:b3:computational-chemistry:split-valence}{diffuse functies} toe.
Reactiebarrières en nauwkeurige energetica: gecorreleerde
golffunctiemethoden op DFT-geometrieën, zover het budget toelaat. Grote
systemen (eiwitten, oppervlakken): een klein gebied dat
kwantummechanisch behandeld wordt binnen een klassiek krachtveld. Test het
gekozen niveau altijd op een verwant molecuul waarvan het antwoord bekend
is.
\end{method}

\begin{inthelab}[Een computationeel experiment]
Een berekening wordt uitgevoerd als een experiment en vastgelegd in het
labjournaal: het programma en zijn versie, de methode, de basis, de
convergentiedrempels, de beginstructuur, en de uitvoerbestanden die als
ruwe gegevens bewaard worden. Een typische reeks is: het molecuul bouwen,
de geometrie optimaliseren op een bescheiden niveau, de frequenties
berekenen om een minimum te bevestigen, en daarna de energie berekenen op
een hoger niveau bij die geometrie.
\end{inthelab}

\begin{history}[Een molecuul dat eerst door berekening bekend was]
% ledger: hist4:HeHp.2019
\ce{HeH+} werd in 1925 in het laboratorium gemaakt. Zijn eigenschappen
werden vanaf de jaren 1930 in detail berekend, en er werd voorspeld dat
het overvloedig voorkomt in sommige astrofysische gassen; maar zijn eerste
rotatielijn ligt in het verre
infrarood en wordt door de atmosfeer geabsorbeerd. In 2019 detecteerde
een vliegend observatorium het in de planetaire nevel NGC 7027, bij de
frequentie die berekeningen en laboratoriumspectra hadden voorspeld.
\end{history}

\section{Oefeningen}

\begin{exercise}[$\star$]\label{exo:b3:computational-chemistry:1}
Neem voor het waterstofatoom de proeffunctie $\eu^{-\alpha r^2}$. In atomaire
eenheden is $E(\alpha) = \frac32\alpha - 2\sqrt{2\alpha/\pi}$. Bepaal de beste
$\alpha$ en de bijbehorende energie; vergelijk met $-\frac12E_h$.
\end{exercise}

\begin{exercise}[$\star$]\label{exo:b3:computational-chemistry:2}
Hoeveel basisfuncties gebruikt een berekening op benzeen in de basis
STO-3G? En in 6-31G(d), met zes cartesische $d$-functies per koolstofatoom
en twee $s$-functies per waterstofatoom?
\end{exercise}

\begin{exercise}[$\star$]\label{exo:b3:computational-chemistry:3}
Leg uit waarom \ce{H2} en \ce{D2} dezelfde bindingslengte en dezelfde
\omterm{def:b3:quantum-model-systems:oscillator}{krachtconstante} hebben, maar verschillende vibratiegolfgetallen. Welke
verhouding verwacht je?
\end{exercise}

\begin{exercise}[$\star$]\label{exo:b3:computational-chemistry:4}
Een frequentieberekening op een structuur met formule \ce{C2H5F} geeft 17
reële frequenties en één imaginaire frequentie, $487\iu$~\unit{cm^{-1}}.
Wat voor \omterm{def:b3:computational-chemistry:pes}{stationair punt} is dit? Hoeveel frequenties verwachtte je in
totaal?
\end{exercise}

\begin{exercise}[$\star\star$]\label{exo:b3:computational-chemistry:5}
% ledger: ie:He, ie4:He2
Bereken met $E(\zeta) = \zeta^2 - \frac{27}{8}\zeta$ de beste
variatie-energie van helium in \unit{\hartree} en in eV, en de fout ervan
ten opzichte van de exacte
$\qty{-2.9034}{\hartree}$. Welke fractie van de totale energie is die fout?
\end{exercise}

\begin{exercise}[$\star\star$]\label{exo:b3:computational-chemistry:6}
Twee basisfuncties geven $H_{11} = \qty{-1.0}{\hartree}$, $H_{22} =
\qty{-0.5}{\hartree}$, $H_{12} = \qty{-0.2}{\hartree}$, $S_{12} = 0.3$. Los de
seculiere vergelijking op. Ligt de laagste wortel onder $H_{11}$?
\end{exercise}

\begin{exercise}[$\star\star$]\label{exo:b3:computational-chemistry:7}
Een hartree-fockberekening op benzeen in STO-3G duurt 10 minuten. Schat de
rekentijd in 6-31G(d), aannemend dat de kost evenredig is met de vierde
macht van het aantal basisfuncties.
\end{exercise}

\begin{exercise}[$\star\star$]\label{exo:b3:computational-chemistry:8}
% ledger: ie4:H2
De STO-3G-orbitaalenergie van \ce{H2} bij \qty{1.4}{\bohr} is
\qty{-0.5782}{\hartree}. Geef de schatting van Koopmans voor de
ionisatie-energie in eV en vergelijk met \qty{15.43}{eV}. Noem twee
redenen waarom ze verschillen, en de richting van elk effect.
\end{exercise}

\begin{exercise}[$\star\star$]\label{exo:b3:computational-chemistry:9}
In STO-3G is de energie van één waterstofatoom \qty{-0.4666}{\hartree} en de
RHF-energie van \ce{H2} bij \qty{10}{\bohr} \qty{-0.5960}{\hartree}. Hoeveel
ligt de RHF-energie op die afstand boven die van twee atomen, in eV?
Verklaar met de ionaire termen.
\end{exercise}

\begin{exercise}[$\star\star\star$]\label{exo:b3:computational-chemistry:10}
Schrijf voor $\phi = c_1\chi_1 + c_2\chi_2$ met reële functies $E(c_1,c_2)$ op
en leid de twee seculiere vergelijkingen af door $\partial E/\partial c_1 =
\partial E/\partial c_2 = 0$ te stellen.
\end{exercise}

\begin{exercise}[$\star\star\star$]\label{exo:b3:computational-chemistry:11}
% ledger: aw:H, diat:H2.we
De RHF/STO-3G-energie van \ce{H2} is $-1.116871$, $-1.117501$ en
$\qty{-1.116714}{\hartree}$ bij $R = 1.30$, 1.35 en \qty{1.40}{\bohr}. Schat
de \omterm{def:b3:quantum-model-systems:oscillator}{krachtconstante} met een eindig verschil, zet haar om in \unit{N/m}
($\qty{1}{\hartree\per\bohr\squared} = \qty{1556.9}{N/m}$), en bereken het
harmonische golfgetal. Vergelijk met $\tilde\omega_e = \qty{4401}{cm^{-1}}$.
\end{exercise}

\begin{exercise}[$\star\star\star$]\label{exo:b3:computational-chemistry:12}
Je moet het energieverschil berekenen tussen twee conformeren van een
suiker, waarvan bekend is dat ze enkele \unit{kJ/mol} verschillen, met
een intramoleculaire waterstofbrug in een van beide. Kies en verantwoord
een methode, een \omterm{def:b3:computational-chemistry:basis}{basisset} en de controles die je zou uitvoeren.
\end{exercise}

\section{Probleem: HeH\texorpdfstring{$^+$}{+}, het eerste molecuul}

\begin{problem}[{Weekendprobleem --- een hartree-fockberekening in de
minimale basis van het heliumhydride-ion: de basis, de core-schatting, de
geconvergeerde orbitaal, en de binding van een proton aan helium}]
\label{pb:b3:computational-chemistry:1}
% ledger: bse:He.sto3g.a1, bse:He.sto3g.a2, bse:He.sto3g.a3, bse:H.sto3g.a1, pa:He, ie:He, ie4:He2
\ce{HeH+} wordt berekend bij $R = \qty{1.4632}{\bohr}$ in de basis STO-3G
(exponenten van He 6.3624, 1.1589, 0.31365; van H 3.4253, 0.62391, 0.16886;
dezelfde drie coëfficiënten). Functie 1 zit op He, functie 2 op H. Het
programma geeft (atomaire eenheden)
\[
\mathbf S = \begin{pmatrix}1 & 0.5368\\ 0.5368 & 1\end{pmatrix},\quad
\mathbf h = \begin{pmatrix}-2.5983 & -1.4318\\ -1.4318 & -1.7318\end{pmatrix},\quad
\mathbf F_{\mathrm{eind}} = \begin{pmatrix}-1.5902 & -1.0610\\ -1.0610 & -0.8340
\end{pmatrix},
\]
de bezette orbitaal $\phi = 0.8766\chi_1 + 0.2025\chi_2$, de orbitaalenergieën
$-1.6328$ en \qty{-0.1725}{\hartree}, en de SCF-energieën $-2.7978$, $-2.8404$,
$-2.8418$, $-2.8418$. Een heliumatoom in dezelfde basis heeft
\qty{-2.8078}{\hartree}. $\qty{1}{\hartree} = \qty{27.211}{eV} =
\qty{2625.5}{kJ/mol}$.

\medskip\noindent\textbf{Deel I --- De basis.}
\begin{enumerate}
  \item Wat betekent ``STO-3G''?
  \item Waarom zijn de exponenten van helium groter dan die van waterstof?
        Controleer dat hun verhouding $(1.69/1.24)^2$ is.
  \item Hoeveel basisfuncties, elektronen en bezette orbitalen zijn er?
  \item Bereken de afstotingsenergie van de kernen.
  \item Interpreteer $S_{12} = 0.5368$.
  \item Waarom is $h_{11}$ lager dan $h_{22}$?
\end{enumerate}

\medskip\noindent\textbf{Deel II --- De core-schatting.}
\begin{enumerate}[resume]
  \item Schrijf de seculiere vergelijking $\det(\mathbf h - \varepsilon\mathbf S) = 0$ als een
        kwadratische vergelijking in $\varepsilon$.
  \item Los haar op.
  \item Waarom is dit slechts een vertrekpunt?
  \item Geef de formule van de fockmatrix uitgedrukt in $\mathbf h$, de
        dichtheidsmatrix en de tweeëlektronintegralen.
  \item Welke wisselwerkingen tussen elektronen bevat $\mathbf F$ die $\mathbf h$ mist?
  \item Leg uit waarom de procedure iteratief moet zijn.
\end{enumerate}

\medskip\noindent\textbf{Deel III --- De geconvergeerde orbitaal.}
\begin{enumerate}[resume]
  \item Controleer dat $\phi$ genormeerd is.
  \item De populaties zijn de diagonaalelementen van $P\mathbf S$, $2c_1^2 + 2c_1c_2S$
        op He en $2c_2^2 + 2c_1c_2S$ op H. Bereken ze.
  \item Waar zit de positieve lading? Wordt het molecuul beter beschreven
        als \ce{He + H+} of als \ce{He+ + H}?
  \item Schat met de \omterm{thm:b3:computational-chemistry:koopmans}{stelling van Koopmans} de energie die nodig is om een
        elektron uit \ce{HeH+} te verwijderen.
  \item Bereken de elektronische energie $E_{\mathrm{el}} = E - V_{\mathrm{nn}}$.
  \item Hoeveel cycli waren nodig om tot op \qty{e-4}{\hartree} te
        convergeren? Waarom daalde de energie bij elke cyclus?
  \item Met exponenten van helium geschaald naar $\zeta = 2.0925$ in plaats
        van 1.69 geeft hetzelfde programma \qty{-2.8607}{\hartree}. Welke
        basis is beter, en waarom kun je dat zeggen?
  \item Wat stelt de energie van de virtuele orbitaal, \qty{-0.1725}{\hartree}, voor?
\end{enumerate}

\medskip\noindent\textbf{Deel IV --- De binding van een proton.}
\begin{enumerate}[resume]
  \item Wat is de energie van een kaal proton? En van \ce{He + H+} ver uit
        elkaar, in deze basis?
  \item Bereken de energie die vrijkomt bij \ce{He + H+ -> HeH+} op dit
        niveau, in \unit{kJ/mol}.
  \item De gemeten protonaffiniteit van helium is \qty{177.8}{kJ/mol}.
        Becommentarieer.
  \item De exacte energie van helium is min de som van zijn
        ionisatie-energieën, 24.5874 en \qty{54.4178}{eV}. Bereken haar in
        \unit{\hartree}, en de fout van het STO-3G-heliumatoom.
  \item Waarom kan een fout van deze grootte toch bruikbare
        molecuulgeometrieën opleveren?
  \item Formuleer het resultaat: de totale RHF/STO-3G-energie van \ce{HeH+} bij
        \qty{1.4632}{\bohr}.
\end{enumerate}
\end{problem}
